% from ega0-20-fr.tex
\oldpage[0\textsubscript{IV-1}]{151}
$B'=B/\mathfrak{N}'$ ; le noyau de l'homomorphisme
$B\otimes_A B\to B'\otimes_{A'}B'$ est
$\mathfrak{U}'=\operatorname{Im}(\mathfrak{N}'\otimes B+B\otimes\mathfrak{N}')$,
d'où résulte aussitôt que le noyau de l'homomorphisme
\begin{equation}
\tag{20.7.14.1}
\label{0.20.7.14.1-fr}
  \varphi_{(\mathfrak{M}',\mathfrak{N}')}:
  \Omega^1_{B/A}\to\Omega^1_{B'/A'}
\end{equation}
est $((\mathfrak{J}\cap\mathfrak{U}')+\mathfrak{J}^2)/\mathfrak{J}^2$ ; par
ailleurs, l'homomorphisme \sref{0.20.7.14.1-fr} est surjectif, comme il
résulte de \sref{0.20.4.7-fr}. Si $\mathfrak{M}''$ (resp.\
$\mathfrak{N}''$) est un second idéal ouvert de $A$ (resp.\ $B$) tel que
$\mathfrak{M}''\subset\mathfrak{M}'$, $\mathfrak{N}''\subset\mathfrak{N}'$
et $\mathfrak{M}''B\subset\mathfrak{N}''$, et si l'on pose
$A''=A/\mathfrak{M}''$, $B''=B/\mathfrak{N}''$, on a de même un
homomorphisme surjectif
\[
  \varphi_{(\mathfrak{M}'',\mathfrak{N}''),(\mathfrak{M}',\mathfrak{N}')}:
  \Omega^1_{B''/A''}\to\Omega^1_{B'/A'}
\]
et ces homomorphismes forment évidemment un \emph{système projectif}. Si
l'on remarque que les
$((\mathfrak{J}\cap\mathfrak{U}')+\mathfrak{J}^2)/\mathfrak{J}^2$ forment
un système fondamental de voisinages de $0$ dans $\Omega^1_{B/A}$, on voit
donc que le \emph{séparé complété} $\widehat{\Omega}^1_{B/A}$ du $B$-module
topologique $\Omega^1_{B/A}$ est donné, à un isomorphisme canonique près, par
\begin{equation}
\tag{20.7.14.2}
\label{0.20.7.14.2-fr}
  \widehat{\Omega}^1_{B/A}=\varprojlim(\Omega^1_{B'/A'}).
\end{equation}

En outre, l'homomorphisme canonique
$j:\Omega^1_{B/A}\to\widehat{\Omega}^1_{B/A}$ est limite projective du
système projectif des $\varphi_{(\mathfrak{M}',\mathfrak{N}')}$, donc
$\varphi_{(\mathfrak{M}',\mathfrak{N}')}$ se factorise en
$\Omega^1_{B/A}\xrightarrow{j}\widehat{\Omega}^1_{B/A}\to
\Omega^1_{B'/A'}$, et comme il est surjectif on en conclut que
l'homomorphisme canonique
\begin{equation}
\tag{20.7.14.3}
\label{0.20.7.14.3-fr}
  \widehat{\Omega}^1_{B/A}\to\Omega^1_{B'/A'}
\end{equation}
est surjectif pour tout couple $(\mathfrak{M}',\mathfrak{N}')$.

On peut d'ailleurs dans ce qui précède remplacer partout $A'$ par $A$.

Enfin, si $L$ est un $B$-module topologique \emph{séparé et complet}, tout
$B$-homomorphisme continu de $\Omega^1_{B/A}$ dans $L$ se prolonge de façon
unique en un $\widehat B$-homomorphisme continu de
$\widehat{\Omega}^1_{B/A}$ dans $L$, et réciproquement un tel homomorphisme
redonne par composition avec $\Omega^1_{B/A}\to\widehat{\Omega}^1_{B/A}$ un
$B$-homomorphisme continu, si bien que l'on a un isomorphisme canonique
\[
  \operatorname{Hom.\ cont}_B(\Omega^1_{B/A},L)\xrightarrow{\sim}
  \operatorname{Hom.\ cont}_{\widehat B}(\widehat{\Omega}^1_{B/A},L).
\]
Plus particulièrement, si $L$ est un $B$-module discret annulé par un idéal
ouvert de $B$, on voit que l'isomorphisme canonique
\sref{0.20.4.8.2-fr} peut aussi s'écrire
\begin{equation}
\tag{20.7.14.4}
\label{0.20.7.14.4-fr}
  \operatorname{Hom.\ cont}_{\widehat B}(\widehat{\Omega}^1_{B/A},L)
  \xrightarrow{\sim}\operatorname{D\acute{e}r.\ cont}_A(B,L).
\end{equation}
\end{env}

\begin{proposition}[20.7.15]
\label{0.20.7.15-fr}
Soient $A$ un anneau discret, $B$ une $A$-algèbre topologique adique
($\mathbf{0}_{\mathrm I}$, 7.1.9), $\mathfrak{n}$ un idéal de définition de
$B$, $B_0=B/\mathfrak{n}$. On suppose que $\Omega^1_{B_0/A}$ et
$\mathfrak{n}/\mathfrak{n}^2$ sont des $B_0$-modules de type fini. Alors
$\widehat{\Omega}^1_{B/A}$ est un $\widehat B$-module de type fini.
\end{proposition}

\begin{proof}
Comme le carré de tout idéal ouvert de $B$ est ouvert, la topologie de
$\Omega^1_{B/A}$ est la topologie $\mathfrak{n}$-préadique
\sref{0.20.4.5-fr}. Compte tenu de l'hypothèse que $B$ est un anneau adique,
il suffit donc, en vertu de ($\mathbf{0}_{\mathrm I}$, 7.2.7 et 7.2.9), de
prouver
